\abstract{
	Niniejsza praca ocenia wydajność zespołowych metod uczenia maszynowego w wykrywaniu i klasyfikacji fałszywych wiadomości na trzech zbiorach danych: \textbf{BuzzFeed}, \textbf{ISOT} oraz \textbf{WELFake}. Zaimplementowano i przetestowano cztery autorskie metody (\texttt{Metoda17}, \texttt{Metoda18}, \texttt{Metoda19}, \texttt{Metoda20}) oraz 16 standardowych metod w trybach uczenia: \textbf{klasycznym}, \textbf{few\-shot} oraz \textbf{one\-shot}, z wykorzystaniem modeli \textbf{BERT}, \textbf{SBERT} i \textbf{RoBERTa}. W trybie \textbf{one-shot}, \texttt{Metoda20} wykazała wyjątkową skuteczność, osiągając niemal idealną dokładność i F1-score (blisko 100\%) dla \textbf{BERT} i \textbf{SBERT} na zbiorach \textbf{BuzzFeed} i \textbf{WELFake}. \texttt{Metoda18} wyróżniła się wysokim AUC-PR (0,8932, 6. miejsce wśród wszystkich metod) na \textbf{WELFake}, demonstrując swoją efektywność w klasyfikacji przy ograniczonych danych. \texttt{Metoda19} osiągnęła znakomite wyniki w trybie \textbf{one-shot} dla \textbf{BERT} na \textbf{BuzzFeed} (dokładność i F1-score bliskie 100\%) oraz solidne rezultaty w trybie \textbf{few-shot} dla \textbf{SBERT} (dokładność 0,7538, F1-score 0,7303). \texttt{Metoda17} pokazała swoją wartość w trybie \textbf{klasycznym}, osiągając AUC-PR 0,7865 (7. miejsce) na \textbf{ISOT}. W trybach \textbf{few-shot} i \textbf{klasycznym}, \texttt{Metoda18} i \texttt{Metoda19} utrzymały wysoką wydajność, z AUC-PR odpowiednio 0,8184 (2. miejsce) na \textbf{ISOT} i 0,8266 na \textbf{WELFake}, przy czasach wykonania poniżej 20 sekund. Analizy statystyczne, w tym testy Shapiro-Wilka, korelacje Pearsona i Spearmana oraz testy Kruskal-Wallis, potwierdziły spójność metryk (korelacja Pearsona >0,99 dla dokładności, MCC i Cohen’s Kappa) oraz zbliżoną wydajność metod (p > 0,05). Zaproponowane autorskie metody charakteryzują się wysoką efektywnością i szybkością obliczeniową, co czyni je pomocniczymi dla wykrywania fałszywych wiadomości, szczególnie w scenariuszach z ograniczonymi danymi.
}{
	This study evaluates the performance of ensemble machine learning methods for detecting and classifying fake news across three datasets: \textbf{BuzzFeed}, \textbf{ISOT}, and \textbf{WELFake}. Four novel methods (\texttt{Metoda17}, \texttt{Metoda18}, \texttt{Metoda19}, \texttt{Metoda20}) and 16 standard methods were implemented and tested in \textbf{classical}, \textbf{few-shot}, and \textbf{one-shot} learning modes, using \textbf{BERT}, \textbf{SBERT}, and \textbf{RoBERTa} models. In the \textbf{one-shot} mode, \texttt{Metoda20} demonstrated exceptional effectiveness, achieving near-perfect accuracy and F1-score (close to 100\%) for \textbf{BERT} and \textbf{SBERT} on \textbf{BuzzFeed} and \textbf{WELFake}. \texttt{Metoda18} stood out with a high AUC-PR (0.8932, 6th place among all methods) on \textbf{WELFake}, showcasing its efficiency in classification with limited data. \texttt{Metoda19} achieved outstanding results in \textbf{one-shot} mode for \textbf{BERT} on \textbf{BuzzFeed} (accuracy and F1-score close to 100\%) and solid performance in \textbf{few-shot} mode for \textbf{SBERT} (accuracy 0.7538, F1-score 0.7303). \texttt{Metoda17} proved its value in \textbf{classical} mode, attaining an AUC-PR of 0.7865 (7th place) on \textbf{ISOT}. In \textbf{few-shot} and \textbf{classical} modes, \texttt{Metoda18} and \texttt{Metoda19} maintained high performance, with AUC-PR scores of 0.8184 (2nd place) on \textbf{ISOT} and 0.8266 on \textbf{WELFake}, respectively, and execution times below 20 seconds. Statistical analyses, including Shapiro-Wilk tests, Pearson and Spearman correlations, and Kruskal-Wallis tests, confirmed metric consistency (Pearson correlation >0.99 for accuracy, MCC, and Cohen’s Kappa) and comparable method performance (p > 0.05). The proposed novel methods exhibit high efficiency and computational speed, making them valuable for fake news detection, particularly in low-data scenarios.
}